\documentclass[instructions.tex]{subfiles}
\begin{document}
 
\chapter{De l'ivrognerie.}
 
\primo Le Sage renvoie le paresseux à la fourmi, pour apprendre à travailler\footnote{Vade ad formicam, ô piger ! Sap.}. Il faut renvoyer l'ivrogne aux bêtes de charge, pour lui apprendre la tempérance. Lorsque l'Église exhorte les pécheurs, elle leur propose l'exemple de Jésus-Christ et des saints; mais il faut changer de langage quand on parle à un intempérant; il faut lui proposer l'exemple des bêtes: oh! quelle horreur!
 
Un ivrogne est indigne de la société des créatures raisonnables. Il faudrait, dit saint Basile, le reléguer parmi les animaux, et le bannir de la compagnie des humains. Les magistrats de la ville de Sparte, dont les habitans étaient les plus sobres de l'univers, ayant exposé en public un esclave plein de vin, pour inspirer l'horreur de ce vice à la jeunesse, le peuple, voyant cet homme dans l'ivresse, fut saisi d'étonnement, et s'écria: Eh! d'où a-t-on fait venir un tel monstre, qui a la figure d'un homme, et qui a moins de sentiment qu'une bête?
 
Où vas-tu, dit un jour un homme de qualité à un de ses subordonnés qui était ivre? Je vais à l'église, répondit-il, prier Dieu. Eh! infâme, répliqua le premier, comment parlerais-tu à Dieu, toi qui n'es pas seulement en état de parler à ton cheval? Prétendre, avec un tel vice, passer pour un homme d'esprit, c'est être dépourvu de bon sens. Telle est l'impudence de l'ivrogne; il ne mérite pas d'être mis au rang des bêtes, et il veut raisonner avec les hommes, et faire le bel esprit\footnote{Cùm ipse insipiens sit, omnes stultos aestimat. Eccles. 10.3.} .
 
\secundo Il n'est point de vice dont on ne puisse corriger un homme avec le secours de la grâce; mais il est très-difficile de corriger un ivrogne. C'est un homme sans foi, sans religion, sans piété, sans respect pour Dieu, pour sa divine parole et pour ses pasteurs. Il se moque des lois de l'Église, des édits du souverain, des règlements des magistrats, qui lui défendent les tavernes.
 
C'est un homme sans pudeur. Non, dit saint Jérôme, je ne croirai jamais qu'un ivrogne soit un homme chaste\footnote{Ebriosum hominm nunquàm castum putabo. S. Hier.}. Il est sans retenue dans ses discours et dans ses chansons, impudent avec les personnes du sexe, dissolu dans ses manières, lubrique dans ses regards, sans respect pour les règles de modestie que la loi de Dieu prescrit dans le mariage.
 
Un intempérant est un homme sans conduite dans ses affaires, sans économie dans sa famille, sans attention sur ses enfans, sans égards, sans charité pour sa femme, pour laquelle il a souvent moins de compassion que pour une bête. Il ne pense, il n'agit, il ne travaille que pour boire; il sacrifie à sa gourmandise ce qu'il doit aux pauvres, ce qu'il doit à ses créanciers et à sa famille. Et, après cent avis qu'on lui a donnés pour le corriger, cet abominable s'étourdit, jusqu'à dire qu'il ne fait aucun mal, qu'il ne fait tort à personne. O aveuglement!
 
III. Pour se convertir, il faudrait qu'il se bornât à une médiocre quantité de vin, qu'il évitât les cabarets, la compagnie des débauchés; qu'il approchât des sacremens, qu'il suivît les avis d'un sage pasteur. Il peut tout cela; il le peut même facilement; mais il ne le fera pas, parce qu'il ne lui plaît pas de le faire, tant son endurcissement est profond.
 
Menacez un intempérant d'une mort funeste, des jugemens de Dieu, des feux de l'enfer, de la perte du ciel; rien ne le touche sur l'état de son âme, dont il a moins de soin que de son chien. Il vit comme s'il n'avait point d'âme; son corps est son maître et son Dieu. Ah! pauvre âme, serais-tu plus déshonorée, si tu étais dans le corps d'un vil animal?
 
On tremble à la vue des malheurs dont Dieu menace les ivrognes dans l'Écriture: Malheur à vous qui êtes forts et puissans à boire! Malheur à vous, qui vous appliquez à vider les pots et les verres! Malheur à vous, qui vous levez le matin pour faire la débauche, qui y passez le jour; pleurez, poussez des cris sur les malheurs qui vous attendent, vous qui faites vos délices du vin! Mais toutes ces menaces d'un Dieu ne font pas plus d'impression sur un ivrogne que sur un rocher.
 
L'Apôtre a eu raison de dire que les ivrognes sont les ennemis de Jésus-Christ; que leur fin sera malheureuse et funeste. Leur aveuglement est si grand, l'état de leur conscience si déplorable, que saint Paul n'en parle qu'en pleurant (1).
 
Va, infâme ivrogne, ce vin que tu bois est comme une couleuvre que tu avales, qui donne la mort à ton âme. Tu n'en crois rien, mais tu apprendras dans les feux éternels qu'il y a un autre Dieu que celui de ton ventre (2).
 
Saint Louis, qui a été l'honneur du trône, fit de sages règlements au sujet des tavernes, avec défense aux cabaretiers, sous de sévères peines, de donner dans leurs logis du vin aux habitans du lieu. Les officiers de justice doivent faire subir aux hôteliers la rigueur des lois, lorsque, contre la défense et les arrêts de police, ils donnent du vin aux personnes dans le lieu de leur domicile. Les confesseurs doivent avoir ici une singulière attention.
 
\section{Résolutions}
 
\primo J'aurai toujours soin d'user de tempérance dans tous mes repas. 

\secundo Je n'y boirai qu'avec modération, presque jamais de vin pur; rarement des liqueurs; tâchant de me priver surtout de quelque chose qui me ferait plaisir, et de sortir de la table avec quelque sentiment d'appétit. 

\prayer{Mon Dieu, faites que je mortifie mon corps, et que je le réduise en servitude, afin qu'il ne devienne pas la cause de ma perte éternelle.}
 
\end{document}
